% !TEX root = ../algebra_02.tex

\section{Классификация простых полей. Простое подполе}

\begin{Definition}{Подполе и простое поле}{}
	Подмножество $F \subset K$ называется подполем, если $F$ является подкольцом в $K$ и само является полем.
	Поле $K$ называется простым, если оно не содержит собственных подполей.
\end{Definition}

\begin{Definition}{Характеристика кольца}{}
	Пусть $A$ --- кольцо. Характеристикой кольца $A$ называется число $n$, определяемое как:
	\[ \operatorname{char} A = \min \{n \in \mathbb{N} \mid \underbrace{1 + \dots + 1}_{n \text{ раз}} = 0\} \]
	Если такого $n$ не существует, полагают $\operatorname{char} A = 0$.
\end{Definition}

\begin{Lemma}{О характеристике поля}{}
	Если $K$ --- поле, то $\operatorname{char} K = 0$ или $\operatorname{char} K = p$, где $p$ --- простое число.
\end{Lemma}

\begin{proof}
	Пусть $\operatorname{char} K = a \cdot b$, где $a > 1, b > 1$. По определению характеристики:
	\[ \underbrace{(1 + \dots + 1)}_{a \text{ раз}} \cdot \underbrace{(1 + \dots + 1)}_{b \text{ раз}} = 0 \]
	Так как в поле нет делителей нуля, хотя бы одна из скобок равна нулю. Это противоречит минимальности характеристики $\operatorname{char} K = ab$. Следовательно, характеристика либо равна нулю, либо является простым числом.
\end{proof}

\begin{Theorem}{О классификации простых полей}{}
	Пусть $K$ --- простое поле. Тогда:
	\begin{enumerate}
		\item Если $\operatorname{char} K = p \neq 0$, то $K \cong \mathbb{F}_p$.
		\item Если $\operatorname{char} K = 0$, то $K \cong \mathbb{Q}$.
	\end{enumerate}
\end{Theorem}

\begin{proof}
	Рассмотрим отображение $\varphi\colon \mathbb{Z} \to K$, заданное правилом:
	\[
	\varphi(n) =
	\begin{cases}
		\underbrace{1 + \dots + 1}_{n \text{ раз}}, & n > 0 \\
		0, & n = 0 \\
		-(\underbrace{1 + \dots + 1}_{-n \text{ раз}}), & n < 0
	\end{cases}
	\]
	Отображение $\varphi$ является гомоморфизмом колец. Ядро этого гомоморфизма зависит от характеристики поля $K$:
	\[
	\ker\varphi =
	\begin{cases}
		(p), & \text{если } \operatorname{char} K = p \neq 0 \\
		0, & \text{если } \operatorname{char} K = 0
	\end{cases}
	\]

	Рассмотрим два случая:
	\begin{enumerate}
		\item Пусть $\operatorname{char} K = p \neq 0$. По теореме о гомоморфизме: $\mathbb{Z}/(p) \cong \operatorname{Im}\varphi $
		
		Поскольку $\mathbb{Z}/(p)$ --- поле, $\operatorname{Im}\varphi$ также является полем и, следовательно, подполем в $K$. Так как $K$ --- простое поле (не содержит собственных подполей), получаем $\operatorname{Im}\varphi = K$.
		Таким образом, $K \cong \mathbb{Z}/(p) = \mathbb{F}_p$.

		\item Пусть $\operatorname{char} K = 0$. Тогда $\ker\varphi = 0$, следовательно, $\varphi$ --- инъективно.
		Построим отображение $\varphi_1\colon \mathbb{Q} \to K$, заданное правилом $\varphi_1\left(\frac{a}{b}\right) = \frac{\varphi(a)}{\varphi(b)}$. (Так как $b \neq 0$, то $\varphi(b) \neq 0$).
		Проверим корректность:
		\[ \frac{a}{b} = \frac{a'}{b'} \implies ab' = a'b \implies \varphi(a)\varphi(b') = \varphi(a')\varphi(b) \implies \frac{\varphi(a)}{\varphi(b)} = \frac{\varphi(a')}{\varphi(b')} \]
		Легко проверить, что $\varphi_1$ является гомоморфизмом колец.
		Найдем ядро: $\varphi_1\left(\frac{a}{b}\right) = 0 \implies \varphi(a) = 0 \implies a = 0 \implies \frac{a}{b} = 0$. Значит, $\ker\varphi_1 = 0$.
		По теореме о гомоморфизме $\mathbb{Q} \cong \operatorname{Im}\varphi_1$.
		Следовательно, $\operatorname{Im}\varphi_1$ --- подполе в $K$. Так как $K$ --- простое поле, $\operatorname{Im}\varphi_1 = K$, откуда $K \cong \mathbb{Q}$.
	\end{enumerate}
\end{proof}

\begin{Proposition}{Простое подполе}{}
	Любое поле содержит простое подполе.
\end{Proposition}

\begin{proof}
	Пусть $K$ --- поле. Рассмотрим пересечение всех его подполей:
	\[ F_0 := \bigcap_{F \text{ --- подполе } K} F \]
	Очевидно, что $F_0$ само является подполем в $K$. Любое подполе $F \subset K$ содержит $F_0$, поэтому $F_0$ не содержит собственных подполей и по определению является простым полем.
\end{proof}
